%! AP Statistics — Unit 2, Lesson 2.8 — Student Packet Cover Sheet
\documentclass[10pt]{article}
\usepackage{apstats-article}
\usepackage{apstats-boxes}
\usepackage{ltablex}
\keepXColumns

\begin{document}

% Full-bleed navy banner
\begin{tikzpicture}[remember picture, overlay]
  \fill[navy] (current page.north west)
    rectangle ([yshift=-0.9in]current page.north east);
\end{tikzpicture}
\vspace{-0.8in}

\begin{center}
  {\color{white}\textbf{\LARGE AP Statistics}}\\[4pt]
  {\color{skymid}\large\textbf{Unit 2: Exploring Two-Variable Data}}\\[6pt]
  {\color{white}\textbf{\large Lesson 2.8 \quad Least Squares Regression}}
\end{center}

\vspace{0.22in}
\namedateperiod
\vspace{0.18in}

% ── LEARNING TARGETS ─────────────────────────────────────────────────────────
\begin{learningtargetbox}
\small
\begin{enumerate}[leftmargin=1.5em, itemsep=5pt]
  \item \ldots{} explain the \textbf{least-squares criterion} and state what quantity the LSRL
        minimizes. \hfill\textcolor{slate}{\small(DAT-1.D)}
  \item \ldots{} compute the \textbf{slope} $b = r\dfrac{s_y}{s_x}$ and
        \textbf{intercept} $a = \bar{y} - b\bar{x}$ from summary statistics.
        \hfill\textcolor{slate}{\small(DAT-1.D)}
  \item \ldots{} compute and interpret the \textbf{coefficient of determination} $r^2$
        using the AP template. \hfill\textcolor{slate}{\small(DAT-1.G)}
  \item \ldots{} identify an \textbf{influential point} and describe its effect on the
        regression line and $r^2$. \hfill\textcolor{slate}{\small(DAT-1.H)}
\end{enumerate}
\end{learningtargetbox}

\vspace{0.14in}

% ── TABLE OF CONTENTS ────────────────────────────────────────────────────────
\begin{tocbox}
\small
\renewcommand{\arraystretch}{1.5}
\begin{tabularx}{\linewidth}{@{} c l X r @{}}
\rowcolor{sky}
\textbf{\#} & \textbf{Component} & \textbf{Description} & \textbf{Score} \\
\midrule
1 & Warm-Up       & Spiral review --- residuals and correlation from Lessons 2.5 \& 2.7  & \blank{1.2cm} \\
2 & Guided Notes  & Least-squares criterion, slope/intercept formulas, $r^2$, influential points & \blank{1.2cm} \\
3 & Group Activity & Compute and interpret the LSRL and $r^2$ from summary statistics        & \blank{1.2cm} \\
4 & Exit Ticket   & Derive and apply the regression line independently                      & \blank{1.2cm} \\
5 & Homework      & Multi-context LSRL, $r^2$, and influential-point practice               & \blank{1.2cm} \\
\midrule
  & \multicolumn{2}{r}{\textbf{Total}} & \blank{1.2cm} \\
\end{tabularx}
\end{tocbox}

\vspace{0.14in}

% ── KEEP IN MIND ─────────────────────────────────────────────────────────────
\begin{remindbox}
\small
The least-squares regression line is the unique line that \textbf{minimizes the sum of squared
residuals}. Two formulas connect it directly to the summary statistics you already know.

\vspace{6pt}
\begin{tabularx}{\linewidth}{@{} >{\centering\arraybackslash\columncolor{goldbg}}p{0.06\linewidth}
                                    X @{}}
\textbf{\large!} &
\textbf{Slope uses $r$, $s_y$, and $s_x$:}\enspace
$b = r\,\dfrac{s_y}{s_x}$.\enspace
The sign of the slope always matches the sign of $r$. \\[6pt]
\textbf{\large!} &
\textbf{The LSRL passes through $(\bar{x},\,\bar{y})$.}\enspace
Use this to find $a$: $\;\;a = \bar{y} - b\bar{x}$. \\[6pt]
\textbf{\large!} &
\textbf{AP $r^2$ template:}\enspace
``About [value]\% of the variation in [response variable] is explained by the
least-squares regression of [response] on [explanatory variable].'' \\[6pt]
\textbf{\large!} &
\textbf{Influential points change the line.}\enspace
A point with an extreme $x$ value that does not follow the overall trend can
dramatically shift the slope, intercept, and $r^2$. Always investigate --- don't
just remove it.
\end{tabularx}
\end{remindbox}

\end{document}
